\section[\nierozwiazane{Zadanie 9}]{Zadanie 9}
 (Zadanie 17.2 z książki) Rozważ następujący wariant sekwencyjnego wrzucania
\(n\) kul do \(n\) urn ponumerowanych liczbami od \(0\) do \(n - 1\). Każda przychodząca kula
losuje jednostajnie liczbę \(i \in \{0, 1, \ldots, n - 1\}\), po czym spośród urn \(i, i + 1, \ldots, i + d - 1 \pmod{n}\)
wybiera tę o najmniejszym ładunku. Udowodnij, że dla ustalonego \(d\) maksymalny
ładunek urny po eksperymencie wynosi \(\Theta(\ln n / \ln \ln n)\).

\textit{Uwaga:} W tym zadaniu można założyć, że \(d = 2\) i \(n\) jest dostatecznie duże.

\brakrozwiazania
